\documentclass{article}

\textwidth=6.5in
\textheight=8.5in
\voffset=-54pt
\oddsidemargin=0in
\evensidemargin=0in
\setlength{\parskip}{8pt}
\setlength{\parindent}{0pt}

\usepackage{workbook}

\usepackage{handout}

\usepackage{exam}

\newcommand{\answerbox}[3][]{
    \fbox{\begin{minipage}[t][#2][c]{#3}
        #1
    \end{minipage}}
    }

\begin{document}

{\large \mbox{} \hfill \bf Name:\underline{\hspace{3in}}}\\[2pt]
{\mbox{} \hfill \it (Please write your name neatly in print.)}\\[12pt]
{\large \mbox{} \hfill \bf HUID:\underline{\hspace{3in}}}
 
\medskip
 
\noindent
{\large\bf
Final Exam\\
Math Mb Spring 2025}
\noindent


{\Large
\begin{center}\textbf{Directions---Please read carefully!} \end{center}}


This exam is subject to the Harvard College Honor Code. No calculators or any other aids are permitted.  Please write as neatly as possible, as answers that can't be read can't be graded and thus can't receive credit. To receive full credit, you must show relevant working
and include units when appropriate. 

You have 180 minutes to complete this exam. 

You may ask for scratch paper if you need additional space to write, but it will not be scanned or graded. Make sure to include all relevant working in this exam booklet in the space provided for each question.

\begin{center}\textbf{Good luck!}\end{center}


\vfill

\centering{
\emph{As a member of Harvard College, I pledge that I will not give or receive assistance on this exam.}}
\vspace{0.3in}

{\large \mbox{} \hfill \bf Signature:\underline{\hspace{3in}}}
\thispagestyle{empty}

\newpage

\begin{enumerate}
    
\item 
\begin{problem}
    The movie \emph{Wicked}, adapted from a musical that was adapted from a novel, is about the story of the Wicked Witch of the West, Elphaba. In the movie, Elphaba escapes from the Wizard of Oz on her flying broomstick and is chased by flying monkeys. 

    As Elphaba flies, she follows the path shown below, which is given by the equation $y^3 + x^3y^2=2x$.

    \begin{center}
    \includegraphics[]{assessments/Final Exam/Elphaba.pdf}
    \end{center}
\end{problem}

    \begin{problem}
        Find an equation of the line tangent to Elphaba's path at the point $(1, 1)$.
    \end{problem}
    \pspace{\vfill}

\pspace{\newpage}

\item 
\begin{problem}
    In the animated movie \emph{Up}, retiree Carl Fredricksen lifts his entire house into the air and flies it to South America using thousands of helium balloons in order to keep a promise to his late wife.
    
    When Carl's house lifts off, a woman standing 60 feet away points her phone camera at the house, rotating her hand to track it while it ascends straight upwards. When her phone's gyroscope measures the angle of elevation to the house to be $\frac{\pi}{6}$ radians, the house is flying upward at 10 feet per second. 
   
   Find the rate at which the woman is rotating her hand when the angle of elevation is $\frac{\pi}{6}$ radians. Give an exact answer in radians per second and fully simplify your answer.

    % \textcolor{red}{
    % (Optional part b involving net change)
    % Before he takes off, Carl's house is 200 feet above sea level. The rate of change of the house's height (in feet above sea level) can be modeled by $r(t)=-\frac{5000t}{t^2+2}$, where $t$ is hours from noon. If Carl flies from 9am to 2pm, what is his elevation above sea level when he lands?
    % }
\end{problem}

% \item 
% \begin{problem}
%     The TV series \emph{Severance} features the mysterious corporation Lumon Industries and its employees. One of the employees, Helly, is trying to escape and is begin chased by the deputy manager Mr. Milchick.

%     Helly has just turned the corner of an L-shaped hallway and is running at a constant speed of 2 m/s. Mr. Milchick is chasing Helly at a constant speed of 3 m/s and is 17 meters away from reaching the corner. Determine how fast the straight-line distance between Helly and Mr. Milchick will be changing \textbf{3 seconds from now.} Show all work and include units in your answer.

%     \textcolor{red}{Should we make the snapshot values easier to find?}
% \end{problem}

\pspace{\newpage}

\item
\begin{problem}
In the movie \emph{The Notebook}, the character Noah Calhoun jumps onto a Ferris wheel while it is moving. The role was played by the actor Ryan Gosling, but jumping onto the Ferris wheel was likely done by an uncredited stunt double.

Suppose the diameter of the Ferris wheel is 40 feet and it makes one counterclockwise revolution every 30 seconds. At time $t=0$, the stunt double jumps onto a seat at the bottom of the wheel, which is 5 feet above the ground.
\end{problem}

\begin{enumerate}
    \item 
    \begin{problem}
        Let $h(t)$ be the stunt double's height above the ground (in feet) at time $t$, where $t$ is measured in seconds. Sketch the graph of $h(t)$ from $t=0$ to $t=90$. Label the axes with appropriate units and clearly indicate the values of the amplitude, midline, and period on the graph.
    \end{problem}
    \pspace{\stretch{2}}

    \item 
    \begin{problem}
        Write a formula for $h(t)$.
    \end{problem}
    \pspace{\stretch{0.5}}

    \item 
    \begin{problem}
        Determine the rate at which the stunt double's height above the ground is changing when $t=5$. Fully simplify your answer and include units.
    \end{problem}
    \pspace{\stretch{2}}

\pspace{\newpage}

    \item 
    \begin{problem}
        Find all times between $t=0$ and $t=30$ when the stunt double is 35 feet above the ground. Fully simplify your answers and include units.
    \end{problem}
    \pspace{\stretch{1}}

    \item 
    \begin{problem}
        Find all times between $t=0$ and $t=30$ when the stunt double is 10 feet above the ground. Give exact answers using inverse trigonometric functions and include units.
    \end{problem}
    \pspace{\stretch{1}}
\end{enumerate}

\pspace{\newpage}

\item 
\begin{problem}
    In the movie \emph{Mean Girls}, the main character Cady Heron wins a math competition by finding a complicated limit. She then goes to her school's Spring Fling dance and is elected Spring Fling queen.
    
    Find the following limits. If the limit does not exist, write DNE, $\infty$, or $-\infty$, as appropriate.

         To receive full credit, you must use correct notation and show all working and/or reasoning. If you use L'Hôpital's Rule, you must check that the rule applies and write $\stackrel{\text{L'H}}{=}$ to indicate where you use it.
\end{problem}

    \begin{enumerate}
        \item 
        \begin{problem}
        $\displaystyle \lim_{x\to 1} \frac{\arctan{x}-\frac{\pi}{4}}{\ln x}$
        \end{problem}
        \pspace{\vfill}
        
        \item 
        \begin{problem}
        $\displaystyle \lim_{x\to0^+} x^{4} \ln{x}$
        \end{problem}
        \pspace{\vfill}

        \pspace{\newpage}
        
        \item 
        \begin{problem}
        $\displaystyle \lim_{x\to \infty} \frac{0.9^x -x^4 + \ln x}{0.5^x + 4x^2 + 1000000}$
        \end{problem}
        \pspace{\vfill}
        
        \item 
        \begin{problem}
        $\displaystyle \lim_{x\to 0^+} \frac{e^{1/x}}{x^2} $
        \end{problem}
        \pspace{\vfill}
    \end{enumerate}

\pspace{\newpage}

\item 
\begin{problem}
    \emph{Stand and Deliver} is a movie dramatization of high school math teacher Jaime Escalante's quest to build a calculus program at Garfield High School in East Los Angeles during the 1970s and 1980s. Though the movie primarily focuses on Escalante's relationships with his students, there are brief scenes that show calculus problems like the ones below.
\end{problem}

\begin{enumerate}
    \item 
    \begin{problem}
        Find $\displaystyle \diff*{x^{1/x}}{x}$.
    \end{problem}
    \pspace{\vfill}

    \item
    \begin{problem}
        Find $\displaystyle \int x^2e^x\,dx$.
    \end{problem}
    \pspace{\vfill}

    \item
    \begin{problem}
        Find $\displaystyle \int \frac{3x}{x+5}\,dx$.
    \end{problem}
    \pspace{\vfill}
\end{enumerate}

\pspace{\newpage}

\item 
\begin{problem}
    In comic books and other media, Barry Allen is the name of a superhero known as ``the Flash'' who has superhuman speed, allowing him to achieve feats like running vertically up the sides of buildings.

    Barry Allen is running up and down the side of a tall skyscraper. His velocity in meters per second is given by the piecewise linear function $v(t)$, where $t$ is measured in seconds and the positive direction is upwards. The graph of $v(t)$ is shown for $0\leq t \leq 8$. At $t=0$, Barry is 25 meters off the ground.
    \vspace{-12pt}
    \begin{center}
        \begin{tikzpicture}
            \begin{axis}[
            xlabel=$t$ (s),
            ylabel=$v(t)$ (m/s),
            xtick distance=1,
            grid=major,
            width=3in
            ]
    \addplot[very thick] coordinates {
      (0,50)
      (2,50)
      (4,-50)
      (7,100)
      (8,50)
    };

    \addplot[closed] coordinates {
      (0,50)
      (2,50)
      (4,-50)
      (7,100)
      (8,50)
    };
            \end{axis}
        \end{tikzpicture}
    \end{center}

    \end{problem}
    
    \begin{enumerate}
        \item 
        \begin{problem}
            State the time interval(s), $0\leq t\leq 8$, during which Barry's height above the ground is \textbf{decreasing}.
        \end{problem}
        \pspace{\stretch{1}}

        \item 
        \begin{enumerate}
            \item 
            \begin{problem}
                State the value(s) of $t$, $0\leq t \leq 8$, when Barry is highest above the ground.
            \end{problem}
            \pspace{\stretch{1}}

            \item
            \begin{problem}
            Find Barry's maximum height above the ground on $0\leq t \leq8$. Include units in your answer.
            \end{problem}
        \end{enumerate}

        \pspace{\stretch{3}}

        \item
        \begin{problem}
            Write, but do not evaluate, an expression involving a definite integral that would give the total distance Barry travels between $t=0$ and $t=8$.
        \end{problem}
        \pspace{\stretch{1}}

\end{enumerate}

\pspace{\newpage}
        
\probonly{
\begin{framed}
        \emph{Continued from the previous page. The problem statement is provided again below.}
        
    Barry Allen is running up and down the side of a tall skyscraper. His velocity in meters per second is given by the piecewise linear function $v(t)$, where $t$ is measured in seconds and the positive direction is upwards. The graph of $v(t)$ is shown for $0\leq t \leq 8$. At $t=0$, Barry is 25 meters off the ground.
    \vspace{-12pt}
    \begin{center}
        \begin{tikzpicture}
            \begin{axis}[
            xlabel=$t$ (s),
            ylabel=$v(t)$ (m/s),
            xtick distance=1,
            grid=major,
            width=3in
            ]
    \addplot[very thick] coordinates {
      (0,50)
      (2,50)
      (4,-50)
      (7,100)
      (8,50)
    };

    \addplot[closed] coordinates {
      (0,50)
      (2,50)
      (4,-50)
      (7,100)
      (8,50)
    };
            \end{axis}
        \end{tikzpicture}
    \end{center}
    \end{framed}
    }

\begin{enumerate}[resume]

        \item 
        \begin{problem}
            Consider the function defined by $\displaystyle f(t) = \int_0^t v(\tau)\,d\tau$. 
            
            Evaluate the expressions below and choose the correct interpretations from the options given.
        \end{problem}

\pspace{6pt}
    \begin{enumerate}
    \item 
    \begin{problem}
        $f(2)=\cfitb{2cm}{100}$
    \end{problem}
    \pspace{12pt}

    \item
    \begin{problem}
    Interpretation of $f(2)$ (circle one): \hfill A \hfill B \hfill C \hfill D \hfill E
    \end{problem}
    \pspace{12pt}

    \item 
    \begin{problem}
        $f'(2)=\cfitb{2cm}{50}$
    \end{problem}
    \pspace{12pt}

    \item
    \begin{problem}
    Interpretation of $f'(2)$ (circle one): \hfill A \hfill B \hfill C \hfill D \hfill E
    \end{problem}
    \pspace{12pt}

      \end{enumerate}

Interpretation options:
      \begin{enumerate}[label=\Alph*.]
      \item Barry's height above the ground (in meters) at $t=2$.
      \item Barry's velocity (in meters per second) at $t=2$.
        \item Barry's acceleration (in meters per second squared) at $t=2$.
      \item The net change in Barry's height above the ground (in meters) from $t=0$ to $t=2$.
      \item The net change in Barry's velocity (in meters per second) from $t=0$ to $t=2$.
      \end{enumerate}
    \end{enumerate}


\pspace{\newpage}

\item 
\begin{problem}
Early computer-animated films like Pixar's \emph{Toy Story} faced production challenges at the intersection of art, computer science, and mathematics. Two aspiring animators are working on the problem below.

Consider the shaded region $\mathcal{R}$ shown. It is bounded above by $y=\arctan x$, on the right by $x=1$, and below by the $x$-axis. Determine the area of $\mathcal{R}$. 
\end{problem}

\hspace{0.5in}
\begin{tikzpicture}
\begin{axis}[
            xlabel=$x$,
            ylabel=$y$,
            ytick=\empty,
            xtick=\empty,
            ymin=-25, ymax=100, width=2.75in]
            
\addplot[thick, name path=F,domain={-1:3}] {atan(x)} ;
\addplot[thick, mark=none] coordinates {(2, -25) (2, 100)};

\addplot[name path=G,domain={0:2}] {0};

\addplot[color=gray!20]fill between[of=F and G, soft clip={domain=0:2}];

\end{axis}
\end{tikzpicture}
\hfill
\begin{tikzpicture}
\begin{axis}[
            xlabel=$x$,
            ylabel=$y$,
            ytick=\empty,
            xtick=\empty,
            ymin=-25, ymax=100, width=2.75in]
            
\addplot[thick, name path=F,domain={-1:3}] {atan(x)} ;
\addplot[thick, mark=none] coordinates {(2, -25) (2, 100)};

\addplot[name path=G,domain={0:2}] {0};

\addplot[color=gray!20]fill between[of=F and G, soft clip={domain=0:2}];

\end{axis}
\end{tikzpicture}
\hspace{0.5in}

Two copies of the same diagram are provided for your convenience. You may annotate the diagrams to illustrate your thinking and communicate your reasoning in parts (a) and (b).

	\begin{enumerate}
	\item
    \begin{problem}
    Andy wants to approximate the area of $\mathcal{R}$ using 100 equal-width \textbf{vertical} rectangles. 

    He needs an expression in sigma notation that he can use a computer to evaluate. Help Andy write an appropriate Riemann sum using sigma notation. Don't use function names like $f$ or $g$ that the computer won't recognize. Make sure to define $\Delta x$ and $x_k$.
    \end{problem}
    
    \probonly{
    \vspace{6pt}
    $\Delta x =$
    \vspace{16pt}
    
    $x_k =$
    \vspace{16pt}

    Riemann sum:
    \vfill
    }

 	\item 

    \begin{problem}
    Barbie wants to approximate the area of $\mathcal{R}$ using 100 equal-width \textbf{horizontal} rectangles. 

    She needs an expression in sigma notation that she can use a computer to evaluate. Help Barbie write an appropriate Riemann sum using sigma notation. Don't use function names like $f$ or $g$ that the computer won't recognize. Make sure to define $\Delta y$ and $y_k$.
    \end{problem}
    
    \probonly{
    \vspace{6pt}
    $\Delta y =$
    \vspace{16pt}
    
    $y_k =$
    \vspace{16pt}

    Riemann sum:
    \vfill
    }


    \pspace{\newpage}



\item
\begin{problem}
    In the limit as the number of rectangles tends to infinity, the Riemann sums from the previous two parts become definite integrals that would give the exact area of $\mathcal{R}$ when evaluated.
\end{problem}

\begin{enumerate}
\item
\begin{problem}
     Write the definite integral that corresponds to Andy's Riemann sum from part (a).
\end{problem}
\pspace{\stretch{0.5}}

\item
\begin{problem}
     Write the definite integral that corresponds to Barbie's Riemann sum from part (b).
\end{problem}
\pspace{\stretch{0.5}}

\item
\begin{problem}
    Choose one of the definite integrals and evaluate it. (\emph{Hint: One may be easier than the other.})
\end{problem}
\pspace{\stretch{3}}
\end{enumerate}
\end{enumerate}

\pspace{\newpage}

\item 
\begin{problem}
    In the movie \emph{Frozen}, the protagonist Elsa has the ability to manipulate ice and snow. In a moment of emotional distress, Elsa releases a magical snowburst of variable density. Snow falls all around her in a perfect circular region with a radius of 9 meters, and the density of the snow is $\rho(r)=\frac{1}{1+r^2}$ kilograms per square meter, where $r$ is the distance to Elsa (in meters).
\end{problem}

\begin{enumerate}
    \item 
    \begin{problem}
        Write, but do not evaluate, a definite integral that would give the total amount of snow created by Elsa's snowburst. 
    \end{problem}
    \pspace{\vfill}

    \item 
    \begin{problem}
        Evaluate the integral you wrote in part (a) to find the total amount of snow created by Elsa's snowburst. Simplify your answer as much as possible and include units.
    \end{problem}
    \pspace{\vfill}
\end{enumerate}
\pspace{\newpage}

\item 
\begin{problem}
In the series \emph{Percy Jackson \& the Olympians}, the protagonists have to navigate a labyrinth. Inside, they encounter a sphinx that asks test questions. In the original Greek myth, the sphinx asks travelers to answer a riddle. If they answer correctly, they are allowed to pass. If they answer incorrectly, they are devoured. You encounter a sphinx that wants you to answer the following calculus question.

What is the derivative of $f(x)=\displaystyle\int_0^1 \sin(t^2)\,dt$?
\end{problem}
\pspace{\vfill}

\item 
\begin{problem}
    In the cartoon \emph{SpongeBob SquarePants,} residents of Bikini Bottom research and catch jellyfish.

    At Jellyfish Rock in Bikini Bottom, a population of jellyfish is growing at a rate proportional to the population size. After doing some research, Sandy discovers that the numerical value of the constant of proportionality is $0.01$ and counts $200$ jellyfish at the time of her research. Let $P(t)$ be the population size at time $t$, where $t$ is the number of hours since Sandy's population count.

    \begin{enumerate}
        
        \item 
        \begin{problem}
        Write down a differential equation that models this situation.
        \end{problem}
        \pspace{\stretch{1}}

        \item
        \begin{problem}
        Write down the general solution to your differential equation from part (a).
        \end{problem}
        \pspace{\stretch{1}}

        \item 
        \begin{problem}
        Find the particular solution that satisfies the initial condition given in the problem.
        \end{problem}
        \pspace{\stretch{2}}

        % \item 
        % \begin{problem}
        %     Find the amount of time it will take for the number of jellyfish at Jellyfish Rock to double.
        % \end{problem}
        % \pspace{\stretch{2}}

\pspace{\newpage}
\end{enumerate}
\probonly{
\begin{framed}
        \emph{Continued from the previous page.}
    \end{framed}
    }

        Jellyfish Caves is home to another colony of jellyfish. Here, the jellyfish population also grows at a rate proportional to the population size, and with the same constant of proportionality, 0.01.
        
        One sunny day, SpongeBob and Squidward decide to go jellyfishing at Jellyfish Caves. The pair catch jellyfish at a constant rate of $6$ jellyfish per hour. Let $J(t)$ be the number of jellyfish at Jellyfish Caves $t$ hours after SpongeBob and Squidward arrive.
        
\begin{enumerate}[resume] 
\item 
\begin{problem}
Which of the following differential equations models this situation? Circle one.
\end{problem}
        \begin{enumerate}[label=\Alph*.]
            \item $\frac{dJ}{dt}=0.01J-6t$
            \item $\frac{dJ}{dt}=0.01J-6$
            \item $\frac{dJ}{dt}=0.01t-6$
            \item $\frac{dJ}{dt}=100e^{0.01t}-6t$
        \end{enumerate}
        \item If there are $100$ jellyfish at Jellyfish Caves when SpongeBob and Squidward arrive, which of the following gives the population of jellyfish in Jellyfish Caves at time $t$? Circle one. To receive full credit, you must show your working and/or reasoning.
        \begin{enumerate}[label=\Alph*.]
            \item $J(t)=(100-6t)e^{0.01t}$
            \item $J(t)=100e^{0.01t}-6t$
            \item $J(t)=600-500e^{0.01t}$
        \end{enumerate}
    \end{enumerate}
\end{problem}

% \probonly{
% \newpage
% \textit{This page has been intentionally left blank for scratch work. No work on this page will be graded.}
% }

% \probonly{
% \newpage
% \textit{This page has been intentionally left blank for scratch work. No work on this page will be graded.}
% }

% \probonly{
% \newpage
% \textit{This page has been intentionally left blank for scratch work. No work on this page will be graded.}
% }

\end{enumerate}

\end{document}*role="textbox", aria-multiline="true", aria-autocomplete="list"*